\section{Experiments}
\label{sec:yield}

\subsection{Experimental setup}
\label{sec:setup}

\paragraph{Dataset.}
We use the \NScans{} scans of LUNA16 with the V2 candidate coordinates and discard their candidate class labels~\cite{setio2017,luna16v2}. After removing one repeated patient and two scans without exactly four reading sessions, the fitting, calibration and test sets contain \NFit{}, \NCal{} and \NTest{} patients, drawn from official subsets \FitSubsets{}, subsets \CalSubsets{}, and subsets \TestSubsets{}, respectively. They contain \NOmitTrain{} training omissions and \NOmitTest{} test omissions among \NTargetTrain{} and \NTargetTest{} targets; matching rules and cohort details are in the supplement.

\paragraph{Detectors and baselines.}
We use two image representations: \NDescriptors{} hand-crafted intensity and Hessian descriptors~\cite{li2003}, and a \EmbedDim{}-dimensional embedding from a 3D CNN applied to \PatchMM\,mm patches, which is trained on reader A's marks from \NDevPatients{} development patients and then frozen. Besides discriminative augmentation, we train heads with first- and second-read labels, nonnegative positive--unlabeled (nnPU) learning, inverse-probability weighting of audited labels (IPW), EM-based and image-only auxiliaries, and the full reference labels as an upper reference. Architectures and optimization details are given in the supplement.

\paragraph{Evaluation protocol.}
The locked comparison uses the C/D reference and $\alpha=\AlphaNom$. On development subsets, we selected half-budget paired-tail sampling with discriminative augmentation and compared it with base-budget mark-stratified IPW. The primary endpoint is recall on test omissions, with a non-inferiority margin of $\NIMargin$ points and additional limits on the alert rate and on false positives per scan (FP/scan). Confidence intervals resample test patients jointly over the six matched audit draws (\NBoot{} bootstrap samples) and condition on the fitted models and the observed calibration patients; calibration-sample variation is assessed in Sec.~\ref{sec:locked}. Single-draw results refer to the first matched draw, the \emph{anchor}. All weights, endpoints and contrasts were fixed before the calibration and test subsets were scored.

\subsection{Audit reach}
\label{sec:reach}

\begin{figure}[t]
\centering
\includegraphics[width=\columnwidth]{acquisition}
\caption{Audit reach: expected number of training omissions with at least one audited candidate (Eq.~\ref{eq:reach}). (a)~All \NOmitTrain{} omissions at the half and base budgets; small points: the six matched draws. (b)~The \NOmitUnmarked{} omissions without any weak mark, at half budget.}
\label{fig:acquisition}
\end{figure}

\begin{figure*}[t]
\centering
\includegraphics[width=\textwidth]{effects}
\caption{Method comparison under two negative pools, with identical scores, targets (\NOmitTest{} omissions, \NTargetTest{} nodules) and \UTest{} common negative test patients. (a)~Mean omission recall \vs mean realized alert rate over calibration-derived thresholds; open markers: nominal $\alpha=\AlphaNom$ (Tab.~\ref{tab:alloc}). (b--d)~Selected method minus each baseline, six-draw mean with \CILevel\% intervals from \NBoot{} bootstrap resamples of the \NTest{} test patients; shaded in (b): beyond the non-inferiority margin.}
\label{fig:effects}
\end{figure*}

On the \NFit{} training patients, half-budget paired-tail sampling audits about \AuditsPair{} candidates and reaches \ReachPair{} of the \NOmitTrain{} omissions in expectation (Fig.~\ref{fig:acquisition}a), whereas mark stratification at the base budget audits about \AuditsMark{} candidates and reaches \ReachMark{}. At the same half budget, image-tail sampling reaches \ReachImage{}, mark stratification \ReachMarkHalf{} and uniform sampling \ReachUniform{}. Most of the gain thus comes from the image ranking, and the full paired-tail policy reaches a further \ReachGain{} omissions on average, between \ReachGainMin{} and \ReachGainMax{} in each of the six matched draws. Realized draws agree, ranging from \DrawPairMin{} to \DrawPairMax{} omissions for paired-tail sampling, \DrawImageMin{} to \DrawImageMax{} for image-tail sampling and \DrawMarkMin{} to \DrawMarkMax{} for base-budget mark stratification. As paired-tail sampling and mark stratification both touch nearly all \NFit{} patients, the gain stems from which candidates are reviewed within each scan rather than from opening more scans.

\paragraph{Sources of the acquisition gain.}
The \ReachGain{}-omission contrast combines the marks with image-dependent uncertainty weighting. We separate them by replaying the saved roster with leverage, budget, floor and cap unchanged (supplement). Image uncertainty without marks ($u_i=p_i$) reaches \ReachImageUnc{} omissions and the marks alone ($p_i=\MarksOnlyP$ in Eq.~\eqref{eq:working}) \ReachMarksOnly{}, so image-dependent uncertainty supplies most of the advantage over image-tail sampling. The complete policy remains the best tested, \AblGain{} omissions or \AblGainPts{} points of reach \AblGainCI{} above image uncertainty alone (unadjusted), but not on the \NOmitUnmarked{} omissions without any weak mark (Fig.~\ref{fig:acquisition}b): there, paired-tail sampling reaches \UnmarkedPair{}, image uncertainty alone \UnmarkedImageUnc{}, image-tail sampling \UnmarkedImage{} and mark stratification \UnmarkedMark{} (\UnmarkedMarkBase{} at the base budget). As $p_i$ predicts the weak marks rather than the reference, this post hoc ablation is practical rather than causal; a factorial over the two leverage terms (supplement) shows that the rank term carries nearly all of the reach.

\subsection{Effect of the representation}

With full reference labels, 3D features recover \OriginalLearned{} of \NTargetTest{} test targets, against \OriginalDescriptor{} for the hand-crafted descriptors, a paired gain of \RepGain{} points \RepGainCI{}; negative-patient alerts fall from \RepAlertsDesc{} to \RepAlertsLearned{} of \PTest{}, and the share of covered targets that outrank every non-target candidate in their scan rises from \RepWinDesc\% to \RepWinLearned\%. As both heads share the same candidate locations, this comparison isolates the representation; we use the 3D features below.

\subsection{Main results}
\label{sec:locked}

\begin{figure*}[t]
\centering
\includegraphics[width=\textwidth]{reference_tail}
\caption{Threshold displacement at fixed scores. (a)~The \NShownMaxima{} largest of the \PCal{} C/D-negative calibration maxima (selected detector); purple diamonds: official consensus nodule. Excluding all \NExcluded{} such patients moves the threshold from rank \PRank{} to the score at rank \PairRank{}, whereas random removal of \NExcluded{} patients keeps it within the shaded ranks \RandomBandLo--\RandomBandHi{} (central \CILevel\% of \NRandom{} draws). (b)~Omissions found among the same \NOmitTest{} targets and (c)~alerts among the same \UTest{} common negative test patients.}
\label{fig:tail}
\end{figure*}

\begin{table}[t]
\centering
\small
\setlength{\tabcolsep}{3.4pt}
\caption{Locked results under the C/D reference at nominal $\alpha=\AlphaNom$, mean of six matched draws. Omit.: recall on \NOmitTest{} test omissions (\%); All: recall on \NTargetTest{} targets (\%); Alert: alerted among the \PTest{} C/D-negative test patients (\%). Half and base are audit budgets; disc.: discriminative augmentation (Eq.~\ref{eq:augmentation}); weighted: tail-weighted objective (supplement). The selected method is in bold.}
\label{tab:alloc}
\begin{tabular}{lrrrr}
\toprule
Method & Omit.\,$\uparrow$ & All\,$\uparrow$ & Alert & FP/scan\,$\downarrow$ \\
\midrule
First read & \textbf{9.5} & \textbf{29.3} & 1.9 & 0.22 \\
nnPU (base) & 7.1 & 22.6 & 1.9 & 0.20 \\
Mark IPW (base) & 7.1 & 20.1 & 4.3 & 0.16 \\
Mark disc. (half) & 7.5 & 24.5 & 4.4 & 0.16 \\
Uniform disc. (half) & 7.9 & 24.9 & 5.9 & 0.16 \\
Image-tail disc. (half) & 4.8 & 23.2 & 2.4 & 0.17 \\
Paired-tail disc. (half)$^\dagger$ & 3.2 & 23.1 & 2.5 & 0.15 \\
Paired-tail IPW (half) & 4.0 & 24.9 & 2.5 & 0.15 \\
Paired-tail weighted (half) & 2.8 & 26.0 & 2.7 & \textbf{0.14} \\
\bottomrule
\end{tabular}

\end{table}

Table~\ref{tab:alloc} reports all training methods. Compared with base-budget mark IPW, the selected method uses \AuditSave\% fewer audits and raises all-target recall from \AllIpw\% to \AllPair\%, a paired gain of \AllGain{} points \AllGainCI{}, at a lower alert rate (\AlertPair\% \vs \AlertIpw\%) and a similar FP/scan (\FPPair{} \vs \FPIpw{}). On the primary endpoint, however, the selected method falls short: its omission recall is \PairPrimary\%, against \IpwPrimary\% for mark IPW, a difference of $\PairRelativeP$ points \PairRelativePCI{} whose lower limit lies below the $\NIMargin$-point margin (Fig.~\ref{fig:effects}b, open circles). The conclusion is unchanged at matched alert rates: selecting each method's threshold on one calibration subset so that their alert rates on the other subset agree yields a difference of $\MatchedDiff$ points \MatchedDiffCI{}. Both intervals condition on the calibration patients; resampling them jointly with the audit draws, with the test patients held fixed, spreads the central \CILevel\% of the omission effect from $\JointLo$ to $+\JointHi$ points (supplement). First-read training, which uses reader A's marks without any audit, attains the highest omission and all-target recall (\OmitFirst\% and \AllFirst\%). With the same estimator and half budget, the mark-stratified, uniform and image-tail allocations reach higher test omission recall (\OmitMarkDisc\%, \OmitUniformDisc\% and \OmitImageDisc\%) than paired-tail allocation, although they audit far fewer training omissions (Sec.~\ref{sec:reach}). Along the full alert curve (Fig.~\ref{fig:effects}a), first-read training leads below an alert rate of \FirstLeadAlert\%, while the audited methods remain close to one another at every rate. The supplement reports a controlled setting in which the first reader keeps only its largest mark per scan (\NOmitBookmark{} imposed test omissions).

\subsection{Analysis of missed omissions}
\label{sec:misses}

\begin{figure*}[t]
\centering
\includegraphics[width=\textwidth]{robustness}
\caption{Robustness of the threshold displacement. (a)~Threshold drop on the original rank scale after excluding official positives (dots) \vs the central \CILevel\% range and median (tick) of \NRandom{} subset-matched random removals. (b)~Omissions found among the same \NOmitTest{} targets when each reader pair's negatives set the threshold; boxed: the original C/D reference.}
\label{fig:robust}
\end{figure*}

Every one of the \NOmitTest{} test omissions is covered by at least one candidate, so the roster itself is not the bottleneck. For the selected detector at the anchor draw, \WinOmit{} omissions outrank every non-target candidate in their scan, yet only one of them exceeds the threshold. Of the \MissedOmit{} missed omissions, \LoseOmit{} are outranked by a non-target candidate in the same scan, whereas the remaining \WinMissed{} win this within-scan comparison but stay below the threshold by a median of \WinGap{} logits; winning within a scan does not, however, establish their ranking against negatives in other scans.

The candidates that determine the threshold are unusual. Among the \NCalNegCand{} candidates on C/D-negative calibration scans, only \NSingleContour{} lie near a contour drawn by just one of the reference readers C and D. Nevertheless, all three sampled maxima around the threshold carry this flag, as do both maxima that trigger alerts on C/D-negative test patients, and both of the latter match official LUNA16 targets. This concentration motivates the negative-pool analysis of Sec.~\ref{sec:tail}.
